\section{Directional-Law Stability and Finite Representation}
This section fixes the measurement class and replaces the directional law. We first
control information and transferred-design performance, then give conditions that
also control optimizer distance and a finite-sample guarantee for an unknown law.
Moment matching gives an algebraic representation guarantee ordered by powers of
SNR, whereas projective-Wasserstein distance gives a finite-SNR metric
perturbation guarantee without requiring polynomial moment matching.

\subsection{Coverage Equivalence and Even Moments}
\begin{definition}[Order-$r$ coverage equivalence]
For an integer $r\ge1$, directional laws $\mu$ and $\nu$ are order-$r$ coverage-equivalent at rank $K$ if
\begin{equation}
M_n(P;\mu)=M_n(P;\nu),\qquad n=1,\ldots,r,
\label{eq:equiv_def}
\end{equation}
for every $P\in\Pset$.
\end{definition}

Every coverage objective is unchanged by replacing $u$ with $-u$, so only the law of $uu^{\top}$ is relevant. Write $U^{\otimes 2r}$ for the $2r$-fold tensor product of $U$ and define the symmetric moment tensor
\begin{equation}
\mathsf T_{2r}(\mu)=\EE_\mu[U^{\otimes 2r}].
\end{equation}
Contracting pairs of indices recovers lower even moments because $\|U\|^2=1$. Tensor equality therefore matches all even polynomial averages of degree at most $2r$ on the sphere.

\begin{lemma}[Rank-one characterization]\label{lem:characterization}
For probability laws $\mu,\nu$ on $\mathbb S^{d-1}$, the following are equivalent:
\begin{enumerate}
\item $\mathsf T_{2r}(\mu)=\mathsf T_{2r}(\nu)$.
\item The laws are order-$r$ coverage-equivalent at rank one.
\item For Gaussian input, the supremum of $|\I_\mu(P;\gamma)-\I_\nu(P;\gamma)|$ over rank-one projectors is $O(\gamma^{r+1})$ as $\gamma\downarrow0$.
\end{enumerate}
These conditions imply order-$r$ coverage equivalence for every rank $K$.
\end{lemma}
\begin{proof}
Tensor equality matches the polynomial $C_P(u)^n$ for every rank and $n\le r$. Conversely, rank-one matching at order $r$ gives equality of $\EE[(v^{\top}U)^{2r}]$ for every unit $v$, hence every $v\in\R^d$ by homogeneity. The coefficients of this polynomial determine the symmetric tensor. Lemma~\ref{lem:expansion} makes moment matching imply the uniform information order. For the converse, each Gaussian Taylor coefficient $(-1)^{n+1}/(2n)$ is nonzero. For every fixed rank-one projector, an information difference of order $O(\gamma^{r+1})$ forces the coefficients through order $r$ to vanish.
\end{proof}

When the target law is Haar measure, matching this tensor is the even-polynomial exactness condition of a weighted real projective $r$-design \cite{hughes2021}. The general object used below is a weighted moment-preserving cubature law; it need not be a classical unweighted design. Odd-polynomial conditions of full spherical designs are unnecessary: the uniform law on $e_1,\ldots,e_d$ and its antipodal symmetrization have identical coverage moments at every order, although the former has nonzero mean. No converse from a single higher rank to tensor equality is needed here.

\subsection{Uniform Information and Design-Transfer Guarantees}
\begin{theorem}[Information-preserving directional representation]\label{thm:representation}
Let $r\ge1$, $\gamma_0>0$, and $h_R\in C^{r+1}([0,\gamma_0])$. The following hold for $0\le\gamma\le\gamma_0$.

\emph{(a) Approximation of objectives and optimal values.} Let
$\varepsilon_n\ge0$, $1\le n\le r$, be bounds on the coverage-moment
errors at a fixed rank $K$, so that
\begin{equation}
\sup_{P\in\Pset}|M_n(P;\mu)-M_n(P;\nu)|\le\varepsilon_n,
\quad 1\le n\le r,
\end{equation}
With the coefficients $b_n$ of Lemma~\ref{lem:expansion}, define
\begin{equation}
\delta(\gamma)=\sum_{n=1}^r|b_n|\gamma^n\varepsilon_n
+\frac{2\gamma^{r+1}}{(r+1)!}\sup_{0\le s\le\gamma}|h_R^{(r+1)}(s)|.
\label{eq:representation_delta}
\end{equation}
Then
\begin{align}
\sup_{P\in\Pset}|\I_\mu(P;\gamma)-\I_\nu(P;\gamma)|&\le\delta(\gamma),\label{eq:equiv_bound}\\
|\I_\mu^\star(\gamma)-\I_\nu^\star(\gamma)|&\le\delta(\gamma).\label{eq:optimal_equiv_bound}
\end{align}
In particular, exact coverage matching sets every $\varepsilon_n$ to zero and gives a uniform $O(\gamma^{r+1})$ error.

\emph{(b) Transfer of measurement designs.} If $P_\nu$ is $\eta$-optimal for the surrogate, meaning $\I_\nu^\star-\I_\nu(P_\nu;\gamma)\le\eta$ for $\eta\ge0$, then
\begin{equation}
0\le\I_\mu^\star(\gamma)-\I_\mu(P_\nu;\gamma)\le2\delta(\gamma)+\eta.
\label{eq:compression_regret}
\end{equation}
The same conclusion holds with any valid uniform objective-error bound replacing $\delta(\gamma)$.

\emph{(c) Finite directional representation.} Write $\delta_u$ for unit
point mass at $u$ and $\operatorname{supp}\mu$ for the smallest closed set
of full $\mu$-measure. For every law $\mu$, there is a weighted law
\begin{equation}
\nu_r=\sum_{\ell=1}^m w_\ell\delta_{u_\ell},\qquad
w_\ell>0,\quad\sum_{\ell=1}^m w_\ell=1,
\end{equation}
with every node $u_\ell$ in $\operatorname{supp}\mu$ and
\begin{equation}
m\le D_{d,r}\triangleq\binom{d+2r-1}{2r},\qquad
\mathsf T_{2r}(\nu_r)=\mathsf T_{2r}(\mu).
\label{eq:support_bound}
\end{equation}
Here $m$ is the number of nodes and $\binom{n}{k}=n!/[k!(n-k)!]$ is the binomial coefficient. The same $\nu_r$ satisfies (a) and (b), with zero moment errors, simultaneously for every measurement rank.
\end{theorem}
\begin{proof}
\emph{(a)} Subtract the uniform expansions in Lemma~\ref{lem:expansion}. Each coefficient difference is bounded by $|b_n|\gamma^n\varepsilon_n$, and the two remainder bounds add. The optimal-value bound follows from $|\max f-\max g|\le\sup|f-g|$.

\emph{(b)} For an original-law optimizer $P_\mu$, insert and subtract the surrogate values at $P_\mu$ and $P_\nu$ in $\I_\mu(P_\mu;\gamma)-\I_\mu(P_\nu;\gamma)$. The two approximation errors contribute at most $2\delta(\gamma)$ and surrogate suboptimality contributes at most $\eta$.

\emph{(c)} Consider the $D_{d,r}$ homogeneous monomials of degree $2r$ restricted to the sphere. They are linearly independent, since a homogeneous polynomial vanishing on the sphere vanishes everywhere by radial scaling. Multiplying a lower even-degree monomial by a power of $\sum_i u_i^2=1$ shows that they span all even polynomials through degree $2r$ on the sphere.

Let $\Phi(u)$ be the vector of these monomials. Its image on the compact set $\operatorname{supp}\mu$ lies in an affine hyperplane, because $(\sum_i u_i^2)^r=1$. The vector $\EE_\mu\Phi(U)$ lies in its convex hull. Carath\'eodory's theorem represents it by at most $D_{d,r}$ image points; discard zero weights. This is the classical finite-dimensional cubature argument underlying Tchakaloff's theorem \cite{bayer2006}. Lemma~\ref{lem:characterization} supplies exact coverage matching for every rank, and (a)--(b) give the information guarantees.
\end{proof}

\begin{corollary}[Gaussian error at every finite SNR]\label{cor:gaussian_representation}
For Gaussian amplitudes and order-$r$ coverage-equivalent laws at rank $K$, let $q=\gamma/(\gamma+2)$ and define
\begin{equation}
\delta_r(\gamma)=\min\left\{\frac{\gamma^{r+1}}{2(r+1)},
\frac{q^{r+1}}{(r+1)(1-q)}\right\}.
\label{eq:compression_delta}
\end{equation}
For every finite $\gamma\ge0$, both the uniform objective error and optimal-value error are at most $\delta_r(\gamma)$. A globally optimal surrogate design has original-law regret at most $2\delta_r(\gamma)$. The finite surrogate of Theorem~\ref{thm:representation}(c) has these guarantees at every rank. At fixed positive finite SNR, the error bound converges geometrically in $r$.
\end{corollary}
\begin{proof}
For $h_R(s)=\tfrac12\log(1+s)$, the exact signed remainder after degree $r$ is
\begin{equation}
\mathcal R_{r+1}^{\rm G}(c)=\frac{(-1)^r}{2}\int_0^{\gamma c}\frac{t^r}{1+t}\,dt.
\label{eq:signed_remainder}
\end{equation}
Its values on $0\le c\le1$ lie in an interval of length at most $\gamma^{r+1}/[2(r+1)]$. After polynomial cancellation, the difference of two expectations is bounded by this interval length.

For the complementary bound, write
\begin{equation}
\tfrac12\log(1+\gamma c)=\tfrac12\log(1+\gamma/2)
+\tfrac12\log(1+q(2c-1)).
\end{equation}
The second logarithm's degree-$r$ polynomial has uniform remainder at most $q^{r+1}/[2(r+1)(1-q)]$. Coverage matching cancels this polynomial and the two remainder bounds add. Theorem~\ref{thm:representation}(a)--(b) transfers the resulting objective bound to optimal values and design regret.
\end{proof}

For a prescribed SNR interval $0\le\gamma\le\Gamma$ with $\Gamma>0$, the same Gaussian surrogate has uniform error at most $\delta_r(\Gamma)$. For $0<\epsilon<1$, set $q_\Gamma=\Gamma/(\Gamma+2)$. It suffices to choose $r\ge1$ with
\begin{equation}
r+1\ge\frac{\log[1/(\epsilon(1-q_\Gamma))]}{-\log q_\Gamma}
\label{eq:error_budget}
\end{equation}
to obtain error at most $\epsilon$ and transfer regret at most $2\epsilon$. At fixed $d$ and $\Gamma$, the sufficient support size is $O((\log(1/\epsilon))^{d-1})$. This is a cubature-existence guarantee, not a computational-efficiency bound; the statistical sample requirement is treated below. For example, $D_{8,1}=36$, $D_{8,2}=330$, and $D_{8,3}=1716$, so the bound need not compress a smaller original law. For finite laws, linear moment constraints provide a constructive route, with residual moment errors covered by Theorem~\ref{thm:representation}(a).

\begin{proposition}[Projective-Wasserstein law perturbation]
\label{prop:wasserstein}
For $u,v\in\mathbb S^{d-1}$, define the projective distance
\begin{equation}
\rho(u,v)=\min\{\|u-v\|,\|u+v\|\}
=\sqrt{2-2|u^{\top}v|}.
\label{eq:projective_distance}
\end{equation}
For laws $\mu$ and $\nu$ on the sphere, define their projective
Wasserstein-1 distance by
\begin{equation}
W_1^{\rm proj}(\mu,\nu)=\inf_{\pi\in\Pi(\mu,\nu)}
\int\rho(u,v)\,d\pi(u,v),
\label{eq:wasserstein_definition}
\end{equation}
where $\Pi(\mu,\nu)$ is the set of couplings with marginals $\mu$ and $\nu$.
Equivalently, this is the Wasserstein-1 distance between the laws induced on
real projective space, so antipodal redistribution has zero distance.
If $h_R\in C^1([0,\gamma])$ and
$L_\gamma=\sup_{0\le s\le\gamma}|h_R'(s)|$, then
\begin{equation}
\sup_{P\in\Pset}|\I_\mu(P;\gamma)-\I_\nu(P;\gamma)|
\le \gamma L_\gamma W_1^{\rm proj}(\mu,\nu).
\label{eq:wasserstein_objective}
\end{equation}
Consequently, every $\eta$-optimal design $P_\nu$ for $\nu$ satisfies
\begin{equation}
\I_\mu^\star(\gamma)-\I_\mu(P_\nu;\gamma)
\le2\gamma L_\gamma W_1^{\rm proj}(\mu,\nu)+\eta.
\label{eq:wasserstein_regret}
\end{equation}
For Gaussian amplitudes, $L_\gamma=1/2$, so the objective error in
\eqref{eq:wasserstein_objective} is at most
$\gamma W_1^{\rm proj}(\mu,\nu)/2$.
\end{proposition}
\begin{proof}
Choose $\varepsilon\in\{-1,1\}$ so that
$\theta=\arccos(u^{\top}\varepsilon v)=\arccos|u^{\top}v|\in[0,\pi/2]$.
The two nonzero eigenvalues of
$uu^{\top}-(\varepsilon v)(\varepsilon v)^{\top}=uu^{\top}-vv^{\top}$ are
$\sin\theta$ and $-\sin\theta$. Ky Fan's inequalities therefore give, for
every $P\in\Pset$,
\begin{equation}
|u^{\top}Pu-v^{\top}Pv|
\le\sin\theta
=\rho(u,v)\cos(\theta/2)
\le\rho(u,v).
\end{equation}
Thus $u\mapsto h_R(\gamma u^{\top}Pu)$ is
$\gamma L_\gamma$-Lipschitz, uniformly in $P$. Integrating under any
coupling and then taking the infimum proves \eqref{eq:wasserstein_objective}.
Theorem~\ref{thm:representation}(b) gives \eqref{eq:wasserstein_regret}.
\end{proof}

Figure~\ref{fig:projective_quotient} shows why the shorter of the two
antipodal chords enters the perturbation bound.
\ProjectiveSphereFigure

\begin{corollary}[Optimizer stability under law replacement]
\label{cor:law_optimizer}
Let
\begin{equation}
e_\gamma=\sup_{P\in\Pset}
|\I_\mu(P;\gamma)-\I_\nu(P;\gamma)|,
\end{equation}
and let $\mathcal M_\mu(\gamma)$ be the nonempty set of maximizers of
$\I_\mu(\cdot;\gamma)$. Define
$\operatorname{dist}_F(P,\mathcal M)=\inf_{Q\in\mathcal M}\|P-Q\|_F$.
Suppose the original-law objective has quadratic growth:
for some $\kappa>0$,
\begin{equation}
\I_\mu^\star(\gamma)-\I_\mu(P;\gamma)
\ge \frac{\kappa}{2}\operatorname{dist}_F
   (P,\mathcal M_\mu(\gamma))^2,
\qquad P\in\Pset.
\label{eq:quadratic_growth_law}
\end{equation}
If $P_\nu$ is $\eta$-optimal for $\nu$, then
\begin{equation}
\operatorname{dist}_F(P_\nu,\mathcal M_\mu(\gamma))
\le \sqrt{\frac{2(2e_\gamma+\eta)}{\kappa}}.
\label{eq:law_optimizer_distance}
\end{equation}
In Theorem~\ref{thm:representation}(a), one may replace $e_\gamma$ by
$\delta(\gamma)$. If the original-law optimizer is unique, this is a direct
Frobenius-distance bound between the transferred and original optimizers.
\end{corollary}
\begin{proof}
The comparison in Theorem~\ref{thm:representation}(b), with the exact uniform
error $e_\gamma$, gives
$\I_\mu^\star-\I_\mu(P_\nu;\gamma)\le2e_\gamma+\eta$.
Combining this inequality with \eqref{eq:quadratic_growth_law} proves
\eqref{eq:law_optimizer_distance}.
\end{proof}

\begin{remark}[When quadratic growth is available]
On the Grassmann manifold, a strict isolated maximizer whose negative objective
has a positive-definite Riemannian Hessian satisfies a local quadratic-growth
condition. Extending that condition to all of $\Pset$, as assumed in
\eqref{eq:quadratic_growth_law}, additionally requires separation of the
objective from its optimum away from the maximizing set. Uniqueness alone does
not imply quadratic growth.
\end{remark}

\begin{corollary}[Joint SNR--directional-law stability]
\label{cor:joint_perturbation}
Let $\gamma,\gamma'\in[0,\Gamma]$, suppose $h_R\in C^1([0,\Gamma])$, and set
\begin{equation}
L_\Gamma=\sup_{0\le s\le\Gamma}|h_R'(s)|,
\qquad
\epsilon_{\gamma,\gamma'}
=L_\Gamma|\gamma-\gamma'|
+L_\Gamma\min\{\gamma,\gamma'\}W_1^{\rm proj}(\mu,\nu).
\label{eq:joint_error}
\end{equation}
Then
\begin{equation}
\sup_{P\in\Pset}|\I_\mu(P;\gamma)-\I_\nu(P;\gamma')|
\le\epsilon_{\gamma,\gamma'}.
\label{eq:joint_objective}
\end{equation}
If $P_{\nu,\gamma'}$ is $\eta$-optimal for law $\nu$ at SNR $\gamma'$, then
\begin{equation}
\I_\mu^\star(\gamma)-\I_\mu(P_{\nu,\gamma'};\gamma)
\le2\epsilon_{\gamma,\gamma'}+\eta.
\label{eq:joint_regret}
\end{equation}
If the target objective $\I_\mu(\cdot;\gamma)$ satisfies
\eqref{eq:quadratic_growth_law} with constant $\kappa$, then
\begin{equation}
\operatorname{dist}_F(P_{\nu,\gamma'},\mathcal M_\mu(\gamma))
\le\sqrt{\frac{2(2\epsilon_{\gamma,\gamma'}+\eta)}{\kappa}}.
\label{eq:joint_optimizer}
\end{equation}
\end{corollary}
\begin{proof}
For every $c\in[0,1]$,
$|h_R(\gamma c)-h_R(\gamma'c)|\le L_\Gamma|\gamma-\gamma'|$.
If $\gamma\le\gamma'$, insert $\I_\nu(P;\gamma)$ between the two objectives;
if $\gamma'\le\gamma$, insert $\I_\mu(P;\gamma')$ instead.
Proposition~\ref{prop:wasserstein} then gives the projective-Wasserstein term at the
smaller SNR, proving \eqref{eq:joint_objective}. The comparison argument of
Theorem~\ref{thm:representation}(b) proves \eqref{eq:joint_regret}, and
Corollary~\ref{cor:law_optimizer} proves \eqref{eq:joint_optimizer}.
\end{proof}

For Gaussian amplitudes, $L_\Gamma=1/2$, so the joint error specializes to
\begin{equation}
\epsilon_{\gamma,\gamma'}
= \frac12|\gamma-\gamma'|
+\frac12\min\{\gamma,\gamma'\}W_1^{\rm proj}(\mu,\nu).
\label{eq:joint_gaussian}
\end{equation}

The next result closes the gap between a known directional law and one
learned from independent direction samples. Its guarantee is uniform over
the continuous projector class, rather than over a finite numerical
candidate set.

\begin{corollary}[Finite-sample learned-law guarantee]
\label{cor:finite_sample}
Let $N,r\ge1$, let $U_1,\ldots,U_N$ be independent samples from $\mu$, let $\gamma>0$, and
suppose $h_R\in C^{r+1}([0,\gamma])$. Define the empirical law
$\widehat\mu_N=N^{-1}\sum_{i=1}^N\delta_{U_i}$ and set
\begin{equation}
L_\gamma=\sup_{0\le s\le\gamma}|h_R'(s)|,
\qquad
H_\gamma=\sup_{0\le s\le\gamma}h_R(s)
-\inf_{0\le s\le\gamma}h_R(s),
\end{equation}
and $p_d=d(d+1)/2$. For $0<\xi\le\sqrt K$ and $0<\alpha<1$, set
\begin{equation}
\Delta_N(\xi,\alpha)=2\gamma L_\gamma\xi
+H_\gamma\sqrt{\frac{p_d\log(1+2\sqrt K/\xi)+\log(2/\alpha)}{2N}}.
\label{eq:empirical_uniform_bound}
\end{equation}
Then, with probability at least $1-\alpha$,
\begin{equation}
\sup_{P\in\Pset}|\I_\mu(P;\gamma)
-\I_{\widehat\mu_N}(P;\gamma)|
\le\Delta_N(\xi,\alpha).
\label{eq:empirical_uniform_information}
\end{equation}
Moreover, there is a weighted law $\nu_{N,r}$ supported on at most
\begin{equation}
\min\left\{N,D_{d,r}\right\}
\end{equation}
observed directions and matching the empirical even moments through degree
$2r$. With
\begin{equation}
\tau_r(\gamma)=\frac{2\gamma^{r+1}}{(r+1)!}
\sup_{0\le s\le\gamma}|h_R^{(r+1)}(s)|,
\end{equation}
the same event satisfies the following bounds for every $\eta$-optimal design
$P_{N,r}$ under $\nu_{N,r}$:
\begin{align}
\sup_{P\in\Pset}|\I_\mu(P;\gamma)-\I_{\nu_{N,r}}(P;\gamma)|
&\le \Delta_N(\xi,\alpha)+\tau_r(\gamma),
\label{eq:sample_compressed_error}\\
\I_\mu^\star(\gamma)-\I_\mu(P_{N,r};\gamma)
&\le2[\Delta_N(\xi,\alpha)+\tau_r(\gamma)]+\eta,
\label{eq:sample_compressed_regret}
\end{align}
For Gaussian amplitudes, $\tau_r(\gamma)$ may be replaced by the sharper
$\delta_r(\gamma)$ in \eqref{eq:compression_delta}. If
\eqref{eq:quadratic_growth_law} holds, the right side of
\eqref{eq:law_optimizer_distance} applies with
$e_\gamma=\Delta_N(\xi,\alpha)+\tau_r(\gamma)$.
\end{corollary}
\begin{proof}
For fixed $u$, the map $P\mapsto h_R(\gamma u^{\top}Pu)$ is
$\gamma L_\gamma$-Lipschitz in Frobenius norm. The rank-$K$ projectors lie in
the radius-$\sqrt K$ ball centered at the origin in the $p_d$-dimensional space of real symmetric
matrices. A maximal $\xi$-separated subset is a $\xi$-net with cardinality at
most $(1+2\sqrt K/\xi)^{p_d}$, by the standard disjoint-ball volume argument.
For every net point, Hoeffding's inequality \cite{hoeffding1963} bounds the
deviation of the empirical average of a variable with range $H_\gamma$.
A union bound over the net gives the square-root term in
\eqref{eq:empirical_uniform_bound}; moving from a projector to its net point
in both the population and empirical objectives adds $2\gamma L_\gamma\xi$.
This proves \eqref{eq:empirical_uniform_information}.

Apply Theorem~\ref{thm:representation}(c) to the finite law
$\widehat\mu_N$. Its nodes can be selected from the observed directions, so
the support contains at most $\min\{N,D_{d,r}\}$ points. Theorem
\ref{thm:representation}(a) bounds the empirical-to-compressed error by
$\tau_r(\gamma)$. The triangle inequality gives
\eqref{eq:sample_compressed_error}, and Theorem~\ref{thm:representation}(b)
gives \eqref{eq:sample_compressed_regret}. The Gaussian and optimizer-distance
claims follow from Corollaries~\ref{cor:gaussian_representation} and
\ref{cor:law_optimizer}, respectively.
\end{proof}

For an explicit accuracy statement, let $L_\gamma>0$, choose a target
$\zeta>0$, and set
$\xi_\zeta=\min\{\sqrt K,\zeta/(4\gamma L_\gamma)\}$. The sufficient condition
\begin{equation}
N\ge \frac{2H_\gamma^2}{\zeta^2}
\left[p_d\log\!\left(1+\frac{2\sqrt K}{\xi_\zeta}\right)
+\log\!\left(\frac2\alpha\right)\right]
\label{eq:sample_complexity}
\end{equation}
gives $\Delta_N(\xi_\zeta,\alpha)\le\zeta$. Thus, apart from the displayed
dimension and logarithmic factors, the directional sample size scales as
$\zeta^{-2}$.

\begin{remark}[Ambient versus intrinsic covering dimension]
The covering argument in \eqref{eq:empirical_uniform_bound} deliberately
embeds $\Pset$ in the $p_d=d(d+1)/2$ dimensional space of symmetric matrices,
which gives an explicit self-contained bound. Direct metric-entropy bounds for
the rank-$K$ Grassmann manifold can replace this ambient exponent by one that
scales with its intrinsic dimension $K(d-K)$, at the cost of additional
covering constants \cite{dai2008}. We retain the ambient argument so the stated
finite-sample constant follows from the elementary volume bound alone.
\end{remark}

\subsection{Spherical Designs and Explicit Ensembles}
A spherical $t$-design reproduces spherical averages of all polynomials of degree at most $t$ \cite{delsarte1977}. It is therefore sufficient, though not necessary, for even-moment matching.
\begin{corollary}[Spherical-design information preservation]\label{cor:design}
Let $\mu_r$ be a spherical $2r$-design and $\sigma$ the Haar law. Under the regularity assumption of Theorem~\ref{thm:representation},
\begin{equation}
\sup_{P\in\Pset}|\I_{\mu_r}(P;\gamma)-\I_\sigma(P;\gamma)|
\le\frac{2\gamma^{r+1}}{(r+1)!}\sup_{0\le s\le\gamma}|h_R^{(r+1)}(s)|.
\label{eq:design_bound}
\end{equation}
For Gaussian input, the sharper bound $\delta_r(\gamma)$ holds.
\end{corollary}
\begin{proof}
Each $C_P(u)^n$ is a polynomial of degree $2n$. A spherical $2r$-design therefore matches coverage moments for $n\le r$. Apply Theorem~\ref{thm:representation}(a) and Corollary~\ref{cor:gaussian_representation}.
\end{proof}
